\par\vspace{5pt}\noindent\begin{minipage}{\linewidth}
\reghead{CORECTL}{0x0E}{R/W (bit 7 per core)}{\texttt{0x01}}
\vspace{2pt}
{\centering\footnotesize\setlength{\tabcolsep}{2pt}%
\begin{tabular}{@{}|>{\centering\arraybackslash}p{12mm}|*{8}{>{\centering\arraybackslash}p{16.9mm}|}@{}}
\hline
\rowcolor{fillD}\textbf{Bit} & \textbf{7} & \textbf{6} & \textbf{5} & \textbf{4} & \textbf{3} & \textbf{2} & \textbf{1} & \textbf{0} \\ \hline
\textbf{Field} & \texttt{ID} & \textcolor{black!45}{--} & \texttt{WAIT1} & \texttt{WAIT0} & \multicolumn{2}{c|}{\textcolor{black!45}{--}} & \texttt{EN1} & \texttt{EN0} \\ \hline
\end{tabular}\par}
\vspace{2pt}
{\small\begin{tabularx}{\linewidth}{@{}llL@{}}
\toprule
Bits & Field & Description \\
\midrule
\texttt{7} & \reg{ID} & Number of the core that reads the register (0 or 1). Read only \\
\texttt{5:4} & \reg{WAIT1, WAIT0} & The core is waiting (1 while executing \texttt{SLEEP}). Read only \\
\texttt{1:0} & \reg{EN1, EN0} & Core enable. Setting a bit starts that core at word 0 of its code memory; clearing it stops the core \\
\bottomrule
\end{tabularx}\par}
\end{minipage}\par

\par\vspace{5pt}\noindent\begin{minipage}{\linewidth}
\reghead{LOCK}{0x05}{R/W (per-core view)}{\texttt{0x00}}
\vspace{2pt}
{\centering\footnotesize\setlength{\tabcolsep}{2pt}%
\begin{tabular}{@{}|>{\centering\arraybackslash}p{12mm}|*{8}{>{\centering\arraybackslash}p{16.9mm}|}@{}}
\hline
\rowcolor{fillD}\textbf{Bit} & \textbf{7} & \textbf{6} & \textbf{5} & \textbf{4} & \textbf{3} & \textbf{2} & \textbf{1} & \textbf{0} \\ \hline
\textbf{Field} & \texttt{OTH3} & \texttt{OTH2} & \texttt{OTH1} & \texttt{OTH0} & \texttt{LK3} & \texttt{LK2} & \texttt{LK1} & \texttt{LK0} \\ \hline
\end{tabular}\par}
\vspace{2pt}
{\small\begin{tabularx}{\linewidth}{@{}llL@{}}
\toprule
Bits & Field & Description \\
\midrule
\texttt{7:4} & \reg{OTH3..OTH0} & Read only: 1 when the other core holds lock $n$ ($n=0..3$) \\
\texttt{3:0} & \reg{LK3..LK0} & Read: 1 when the reading core holds lock $n$. Write 1 to take the lock if free; write 0 to release it if held by the writing core \\
\bottomrule
\end{tabularx}\par}
\end{minipage}\par

\par\vspace{5pt}\noindent\begin{minipage}{\linewidth}
\reghead{ARB}{0x3F50}{R/W, protected}{\texttt{0x00}}
\vspace{2pt}
{\centering\footnotesize\setlength{\tabcolsep}{2pt}%
\begin{tabular}{@{}|>{\centering\arraybackslash}p{12mm}|*{8}{>{\centering\arraybackslash}p{16.9mm}|}@{}}
\hline
\rowcolor{fillD}\textbf{Bit} & \textbf{7} & \textbf{6} & \textbf{5} & \textbf{4} & \textbf{3} & \textbf{2} & \textbf{1} & \textbf{0} \\ \hline
\textbf{Field} & \multicolumn{6}{c|}{\textcolor{black!45}{--}} & \multicolumn{2}{c|}{\texttt{MODE}} \\ \hline
\end{tabular}\par}
\vspace{2pt}
{\small\begin{tabularx}{\linewidth}{@{}llL@{}}
\toprule
Bits & Field & Description \\
\midrule
\texttt{1:0} & \reg{MODE} & Bus arbitration when both cores request the shared bus: 0 alternating (first collision after reset goes to core 0), 1 always core 0, 2 always core 1, 3 reserved \\
\bottomrule
\end{tabularx}\par}
\end{minipage}\par